\begin{textsample}{POS Dim 1 – 2017\_c4 – Score -3.00 – 2017\_c4/enade\_2017\_eng\_civ\_9cc7e908b9.txt}  \label{ex:f1_pos_277}
CAPÍTULO 4 PERCEPÇÃO DA PROVA As análises feitas neste capítulo tratam das percepções dos concluintes da Área de Engenharia Civil sobre a prova aplicada no Enade/2017.
Estas percepções foram mensuradas por meio de nove questões que avaliaram desde o grau de dificuldade da prova até o tempo gasto para concluí -la.
As percepções sobre a prova foram relacionadas com o desempenho dos estudantes e com a Grande Região de funcionamento do curso.
O questionário de percepção da prova encontra -se ao final do Anexo VII, que traz a reprodução da prova.
O desempenho dos estudantes foi classificado em quatro quartos.
Para tanto, esse desempenho foi ordenado de forma ascendente.
O percentil 25, P25, também conhecido como primeiro quartil, é a nota de desempenho que deixa um quarto ( 25 \% ) dos valores observados abaixo e três quartos acima.
A Figura 1 apresenta uma ilustração deste conceito.
O quarto inferior de desempenho é composto pelas notas abaixo do primeiro quartil.
Já o percentil 75, P75, também conhecido como terceiro quartil, é o valor para o qual há três quartos ( 75 \% ) dos dados abaixo e um quarto acima dele.
O quarto superior de desempenho é composto pelas notas iguais ou acima do terceiro quartil.
O percentil 50, P50, também conhecido como mediana, é o valor que divide as notas em dois conjuntos de igual tamanho.
O segundo quarto inclui valores entre o primeiro quartil ( P25 ) e a mediana.
O terceiro quarto contém os valores entre a mediana ( P50 ) e o terceiro quartil ( P75 ).
Vale ressaltar que percentis, quartis e medianas são pontos que não obrigatoriamente pertencem ao conjunto original de dados, ao passo que os quartos são subconjuntos dos dados originais.
A seguir, serão apresentados gráficos com resultados selecionados, relativos às nove questões avaliadas por grupos de estudantes.
Os gráficos apresentam nas barras o percentual de alunos que assinalaram uma das opções ou a soma das porcentagens daqueles que assinalaram duas ( ou três ) delas.
Por exemplo, para as questões 1 e 2, os gráficos apresentam a porcentagem total de participantes que assinalaram as opções ( D ) Difícil e ( E ) Muito difícil.
Em cada barra foram assinalados também os extremos do intervalo de confiança de 95 \% como linhas verticais unidas por uma linha horizontal na forma da letra H maiúscula.
As Tabelas no Anexo II apresentam os valores absolutos e a distribuição percentual17 das alternativas válidas das nove questões, segundo o mesmo recorte de desempenho dos alunos e Grande Região de funcionamento do curso. 4.1 GRAU DE DIFICULDADE DA PROVA 4.1.1 Componente de Formação Geral Ao avaliarem Qual o grau de dificuldade desta prova na parte de Formação Geral? ( Questão 1 ), 27,0 \% do grupo de inscritos e presentes optaram pelas alternativas Difícil ou Muito difícil.
Entretanto, para mais da metade dos estudantes ( 60,7 \% ), o Componente de Formação Geral da prova foi considerado com grau de dificuldade Médio ( Gráfico 4.1, Gráfico 4.2 e, no Anexo II, a Tabela II.1 ).
O percentual de estudantes que consideraram a prova como Difícil ou Muito difícil foi maior na região Sudeste, onde a proporção foi de 29,0 \%, enquanto a de menor incidência foi a Nordeste, com 21,4 \%.
No Gráfico 4.1, é possível observar que as diferenças entre a região Nordeste e as demais regiões são estatisticamente significativas, bem como as diferenças entre o Sudeste e as demais regiões.
Nas Grandes Regiões, a proporção de presentes à prova que consideraram o Componente de Formação Geral como sendo de grau de dificuldade Médio esteve entre 59,4 \% na região Sudeste, e 64,4 \% na região Nordeste. 17 Cumpre lembrar uma das convenções para tabelas numéricas ( pág. iii ) sobre a possibilidade de a soma das partes não resultar em 100 \% por questões de arredondamento.
O percentual de alunos que consideraram a prova Difícil ou Muito difícil foi decrescente em função dos quartos de desempenho : 34,9 \% no primeiro quarto e 17,9 \% no quarto, grupo de melhor desempenho na prova.
Nos quartos de desempenho intermediários, a proporção de alunos que consideraram a prova Difícil ou Muito difícil foi de 30,0 \% no segundo quarto e 25,1 \% no terceiro quarto.
As diferenças entre os quartos são estatisticamente significativas.
Para todos os quartos de desempenho, a alternativa modal para essa pergunta foi Médio, com 56,9 \% e 63,0 \% dos respondentes nos quartos extremos, primeiro e quarto, respectivamente. 4.1.2 Componente de Conhecimento Específico Ao responderem à Questão 2 – Qual o grau de dificuldade desta prova na parte de Componente Específico? – 50,8 \% do grupo de estudantes classificaram -na como Difícil ou Muito difícil.
Além disso, o Componente de Conhecimento Específico da prova foi considerado com grau de dificuldade Médio por 45,7 \% dos alunos ( Gráfico 4.3, Gráfico 4.4, e, no Anexo II, a Tabela II.2 ).
A análise das respostas dos estudantes quanto ao grau de dificuldade do Componente de Conhecimento Específico da prova, agregado por Grande Região, mostra que a diferença entre a maior e as duas menores proporções de alunos que a avaliaram como Difícil ou Muito difícil é estatisticamente significativa : a maior na região Sul ( 53,6 \% ) e as duas menores proporções nas regiões Nordeste ( 44,1 \% ) e Norte ( 49,5 \% ).
O percentual de alunos que classificaram o grau de dificuldade como Médio, no Componente de Conhecimento Específico, variou de 43,5 \% a 51,9 \%, para as regiões Sul e Nordeste, respectivamente.
Considerando -se a avaliação da dificuldade das questões do Componente de Conhecimento Específico da prova, de acordo com o desempenho dos estudantes, observa -se que há diferença estatisticamente significativa dos resultados entre o quarto de desempenho inferior e os demais, bem como entre o quarto superior e os demais.
A proporção dos que classificaram a parte específica como Difícil ou Muito difícil variou de 45,3 \% ( último quarto ) a 54,5 \% ( segundo quarto ).
A proporção dos que classificaram a Questão 2 como de grau Médio foi de 44,9 \% do quarto inferior e 50,3 \% do superior. 4.2 EXTENSÃO DA PROVA EM RELAÇÃO AO TEMPO TOTAL Indagados quanto à extensão da prova, em relação ao tempo total oferecido para a sua resolução ( Questão 3 ), os estudantes apontaram, com maior incidência, a alternativa que considerava a extensão adequada, para todas as agregações consideradas ( Gráfico 4.5, Gráfico 4.6, e, no Anexo II, a Tabela II.3 ).
O percentual de alunos que responderam ser a extensão da prova adequada foi de 53,8 \%.
Já 37,3 \% dos inscritos presentes consideraram que a prova foi longa ou muito longa, e 8,9 \% a avaliaram como curta ou muito curta.
Entre as Grandes Regiões, a proporção daqueles que avaliaram a prova como longa ou muito longa em relação ao tempo total destinado à sua resolução variou de 32,5 \% na região Norte até 39,2 \% na região Nordeste.
A diferença entre a região Norte e as demais ( à exceção da Centro-Oeste ) é estatisticamente significativa, bem como a diferença da região Nordeste para as demais ( à exceção da Sul ). exceção da Centro-Oeste ) é estatisticamente significativa, bem como a diferença da região Nordeste para as demais ( à exceção da Sul ).
Considerando -se o desempenho dos alunos, nota -se ainda que 55,3 \% consideraram a extensão da prova adequada no quarto de desempenho inferior, e 53,5 \%, no de melhor desempenho ( quarto superior ).
Nos quartos intermediários, essa proporção foi de 53,5 \% no segundo quarto e de 52,9 \% no terceiro.
No Gráfico 4.6, pode -se constatar que há uma tendência crescente da proporção de estudantes que consideraram a prova longa ou muito longa em função dos quartos de desempenho.
Observa -se diferença estatisticamente significativa entre as proporções dos quartos inferior e superior. 4.3 COMPREENSÃO DOS ENUNCIADOS DAS QUESTÕES 4.3.1 Componente de Formação Geral Em relação aos enunciados das questões do Componente de Formação Geral ( Questão 4 ), as opiniões foram positivas, já que 76,6 \% dos alunos avaliados consideraram todos ou a maioria dos enunciados das questões claros e objetivos ( Gráfico 4.7, Gráfico 4.8, e, no Anexo II, a Tabela II.4 ).
Na análise regional, a percentagem de estudantes que avaliaram que todos ou a maioria dos enunciados das questões do Componente de Formação Geral estavam claros e objetivos variou de 76,3 \% na região Sudeste a 78,2 \% na região Centro-Oeste, não ocorrendo diferenças estatisticamente significativas entre as regiões.
A análise das percepções dos estudantes sobre a clareza e objetividade dos enunciados permite afirmar que todos, ou a maioria dos enunciados de questões relativas ao Componente de Formação Geral, foram considerados claros e objetivos para a maior parte dos respondentes ( maior ou igual a 76,3 \% em todas as regiões e maior ou igual a 70,2 \% para todos os quartos de desempenho ).
Segundo o desempenho, observa -se que a proporção dos que emitiram essa opinião cresce conforme o desempenho aumenta, com diferenças estatisticamente significativas entre todos os quartos de desempenho.
No quarto superior, a clareza e objetividade de todos ou da maioria dos enunciados das questões foi percebida por 82,2 \%. 4.3.2 Componente de Conhecimento Específico Em relação aos enunciados das questões do Componente de Conhecimento Específico da prova, para 74,7 \% dos estudantes avaliados da Área de Engenharia Civil, a clareza e a objetividade ( Questão 5 ) estavam presentes em todas ou na maioria das questões ( Gráfico 4.9, Gráfico 4.10, e no Anexo II, a Tabela II.5 ).
A maioria dos estudantes de todas as Grandes Regiões brasileiras considerou claros e objetivos todos ou a maioria dos enunciados das questões do Componente de Conhecimento Específico da prova, percentual sempre maior ou igual a 73,3 \%.
As diferenças da região Sudeste em relação às demais ( à exceção da Norte ) são estatisticamente significativas, bem como as diferenças entre as regiões Centro-Oeste e Norte.
A proporção de estudantes que consideraram os enunciados das questões claros e objetivos apresenta uma tendência crescente em relação ao aumento de desempenho : mais elevada no quarto superior ( 82,3 \% ) se comparada ao quarto inferior de desempenho ( 67,9 \% ).
As diferenças entre todos os quartos de desempenho são estatisticamente significativas. 4.4 SUFICIÊNCIA DAS INFORMAÇÕES/INSTRUÇÕES FORNECIDAS Ao se avaliarem as informações/instruções fornecidas para a resolução das questões ( Questão 6 ), 82,3 \% dos respondentes da Área de Engenharia Civil de todo o Brasil afirmaram que estas eram até excessivas ou suficientes em todas ou na maioria das questões ( Gráfico 4.11, Gráfico 4.12, e, no Anexo II, a Tabela II.6 ).
Quanto à distribuição de respondentes pelas Grandes Regiões, observa -se que a proporção de estudantes que consideraram as informações/instruções fornecidas até excessivas ou suficientes em todas ou na maioria das questões foi sempre superior ou igual a 80,7 \%, chegando a 85,6 \% na região Sul.
A diferença entre o maior percentual, obtido na região Sul, e o menor, da região Norte ( 80,7 \% ) é estatisticamente significativa.
Levando -se em conta o desempenho dos participantes, notam -se diferenças estatisticamente significativas entre as opiniões de estudantes de todos os quartos de desempenho, como mostra o Gráfico 4.12.
O percentual de participantes que avaliaram as informações/instruções como até excessivas ou suficientes em todas ou na maioria das questões foi mais elevado no quarto superior ( 89,5 \% ), percentual superior à média nacional ( 82,3 \% ).
Já no quarto inferior, a suficiência das informações/instruções declarada como até excessiva, em todas ou na maioria das questões foi percebida por 75,2 \% dos respondentes.
Essas proporções são crescentes com os quartos de desempenho. 4.5 DIFICULDADE ENCONTRADA AO RESPONDER À PROVA Indagados sobre as dificuldades com as quais se depararam ao responder à prova ( Questão 7 ), 17,5 \% dos estudantes apontaram o Desconhecimento do conteúdo.
Para 51,2 \%, a Forma diferente de abordagem do conteúdo foi indicada como dificuldade.
Já a Falta de motivação para fazer a prova foi a dificuldade apontada por 12,0 \% dos respondentes.
Considerando -se todo o Brasil, 11,8 \% dos respondentes afirmaram que não tiveram qualquer tipo de dificuldade para responder à prova ( Tabela II.7 no Anexo II ).
Os Gráficos 4.13 e 4.14 apresentam os percentuais de estudantes que apontaram o Desconhecimento do conteúdo como dificuldade percebida ao responder à prova.
Na análise por Grandes Regiões, o percentual de inscritos e presentes que apontaram o Desconhecimento do conteúdo como dificuldade ao responder à prova não superou 19,9 \%.
Os percentuais variaram de 16,3 \% na região Nordeste a 19,9 \% na região Centro-Oeste, sendo esta diferença significativa estatisticamente.
A Forma diferente de abordagem do conteúdo foi a escolha modal dos estudantes, com percentuais que variaram de 46,3 \% ( região Nordeste ) a 52,9 \% ( região Sudeste ).
O percentual de alunos que citaram a Falta de motivação como dificuldade variou de 10,4 \% ( região Sudeste ) a 14,7 \% ( região Nordeste ).
Os que declararam não ter qualquer tipo de dificuldade para responder à prova variaram de 10,3 \% nas regiões Norte e Centro-Oeste a 12,9 \% na região Nordeste.
Com relação aos quartos de desempenho, o Desconhecimento do conteúdo foi a opção escolhida por 18,6 \% dos estudantes do quarto superior e por 17,3 \% daqueles do quarto inferior.
Nota -se que a proporção do segundo quarto apresenta diferença estatisticamente significativa para o quarto superior.
A alternativa modal para os alunos, quando agregados pelos quartos de desempenho, foi que a dificuldade encontrada foi causada pela Forma diferente de abordagem do conteúdo : 50,5 \% no quarto inferior e 45,4 \% do quarto superior assim o responderam. 4.6 CONTEÚDOS DAS QUESTÕES OBJETIVAS DA PROVA Ao se analisarem os conteúdos das questões objetivas da prova ( Questão 8 ), verifica -se que um percentual muito pequeno dos estudantes avaliados, apenas 3,6 \%, afirmou que não estudou ainda a maioria desses conteúdos ( Gráficos 4.15, Gráfico 4.16, e a Tabelas II.8 no Anexo II ).
A maioria ( 67,3 \% ) afirmou ter estudado e aprendido muitos ou todos os conteúdos.
Na análise por Grande Região, a proporção de respondentes que escolheram a opção não estudou ainda a maioria desses conteúdos foi pequena.
Observa -se que apenas na região Sul ( 2,2 \% ) a proporção é menor que a média nacional ( 3,6 \% ).
Observa -se diferença estatisticamente significativa entre a região Sul e as demais.
Em todas as regiões, a maioria dos presentes afirmou ter estudado e aprendido muitos ou todos os conteúdos, com proporções variando entre 61,7 \% na região Norte e 71,4 \% na Sul.
Considerando -se separadamente as opiniões de estudantes dos quatro quartos de desempenho, observa -se que, no quarto inferior, 6,6 \% ofereceram como resposta que não estudou ainda a maioria desses conteúdos, sendo 1,5 \% os do quarto superior com a mesma resposta.
Só não se observa diferença estatisticamente significativa entre os alunos que optaram por este motivo de dificuldade entre os quartos intermediários de desempenho ( segundo e terceiro ).
Tendo em conta o quarto superior, 79,1 \% dos alunos afirmaram ter estudado e aprendido muitos ou todos os conteúdos.
No outro extremo, no primeiro quarto, 55,9 \% optaram pelas mesmas categorias.
As proporções são crescentes com o desempenho. 4.7 TEMPO GASTO PARA CONCLUIR A PROVA Ao responderem sobre o tempo de conclusão da prova ( Questão 9 ), quase três quartos dos estudantes ( 76,5 \% ) afirmaram ter gasto Entre duas e quatro horas ( Gráfico 4.17, Gráfico 4.18 e, no Anexo II, a Tabela II.9 ).
Considerando -se as cinco Grandes Regiões brasileiras, os percentuais dos que utilizaram Entre duas e quatro horas para finalizar a prova nas regiões Nordeste ( 72,9 \% ), Centro-Oeste ( 74,1 \% ) e Norte ( 74,7 \% ) foram inferiores ao percentual nacional.
Nas demais Grandes Regiões, o percentual de alunos que dispensaram Entre duas e quatro horas para concluir a prova ficou igual ou acima de 77,3 \%, como mostra o Gráfico 4.17.
Há diferenças estatisticamente significativas entre várias regiões.
Uma vez sendo analisadas as alternativas escolhidas pelos estudantes que se situam nos diferentes quartos de desempenho, observa -se uma tendência crescente : uma proporção maior de participantes no quarto superior ( 80,1 \% ) declarou ter gasto Entre duas e quatro horas para concluir a prova quando comparada à proporção dos quartos inferiores.
As diferenças entre o primeiro ( 70,5 \% ) e os demais quartos de desempenho são estatisticamente significativas, bem como entre o quarto superior e os demais quartos de desempenho.

% matched lemmas: 
\end{textsample}
